\section{Summary}
\label{sec:conclusion:summary}

This Phase~1 report has presented the design of \textbf{BusLens},
a visual analytics platform for the Ho Chi Minh City public bus
network. Chapter~\ref{chapter:intro} motivated the work as the
engineering response to the absence of any analytical tool over
the network's three large data streams (raw GPS pings, AFC
ticketing events, and static GTFS metadata).
Chapter~\ref{chapter:background} surveyed four bodies of prior
work --- bus operational analytics (Srisakaokul, Derevitskiy,
Lin~\&~Zeng), visualisation of spatial-temporal traffic data
(Liu), bus--metro intermodal connectivity, and existing HCMC
tooling --- and developed the theoretical foundations of the
solution: the four-layer warehouse adapted from the medallion
pattern~\cite{databricksMedallion}, the analytical-stack choice
(DuckDB~\cite{duckdb} + Streamlit~\cite{streamlit} +
Plotly~\cite{plotly} + pydeck~\cite{pydeck} +
\texttt{st.components.v1.html}), the QoS scoring
framework~\cite{srisakaokul2023}, the visualisation primitives
of Liu~\cite{liu2022} plus the Marey time--distance
diagram~\cite{linzeng2014}, the visual-storytelling principles
of Knaflic~\cite{knaflic2015}, and the Multiple Coordinated
Views interaction model that the dashboard's linked-brushing
invariant implements.

Chapter~\ref{chapter:proposed} developed the proposed solution
along six axes: a formal problem statement; an SDVF tier
prioritisation; a four-layer warehouse design with eight
aggregate tables that compose a hard data contract with the
chart catalogue; the EDA findings that scope every downstream
decision; a thirty-one chart catalogue with chart-to-aggregate
mapping and SDVF tiers; and the dashboard's
Hub~$\rightarrow$~Card~$\rightarrow$~Chart information
architecture with the A1, A2, A3 behavioural invariants and the
Stories scrollytelling layer.
Chapter~\ref{chapter:prototype} documented the four executable
deliverables of Phase~1: the EDA notebooks, the UI mockup, the
warehouse skeleton, and the supporting design artefacts.
Chapter~\ref{chapter:future} sketched the Phase~2
implementation roadmap, the data activations that unlock the
deferred charts, and the predictive extensions and external
integrations that follow the historical foundation.

\section{Achievements}
\label{sec:conclusion:achievements}

Three achievements of this Phase~1 work deserve emphasis.

\paragraph{A unified data foundation.}
The four-layer warehouse design reconciles three previously
siloed data streams behind a single SQL surface, with
reproducible normalisation from the upstream parquet exports
and consistent type-casting in
\texttt{src/duckdb\_loader.py}. The foundation is reusable for
any future analytical work on the same data, including the
natural-language ``Ask the Data'' interface under development
by another team
(Section~\ref{sec:future:external}).

\paragraph{A research-grounded analytics suite.}
Each visualisation primitive used in the dashboard traces back
to a specific peer-reviewed source. The QoS-1/2/3 scoring
framework of Srisakaokul~\textit{et~al.}~\cite{srisakaokul2023},
originally delivered as an offline study, is wired into an
interactive tool. The FuzzyCut, GroupSet, and Gridify primitives
of Liu~\cite{liu2022}, originally demonstrated on highway-sensor
and taxi data, are transferred to the urban-bus operational
domain. The Marey time--distance
diagram~\cite{linzeng2014} is the principal Detail-mode chart
of the Routes card. The MCV interaction model is implemented at
full scope across the dashboard as the A1 invariant. The
result is a tool whose every analytical claim is grounded in a
citable source.

\paragraph{A unified, role-anchored interactive interface.}
The Hub~$\rightarrow$~Card~$\rightarrow$~Chart hierarchy, the
single-chart canvas rule, the three A1/A2/A3 behavioural
invariants, and the shared session-state object together
deliver a dashboard specification that scales from a
1366$\times$768 laptop screen to a defence-room projector and
serves four distinct user roles (passenger, driver, company,
MCPT) from a single chrome without forcing any of them to leave
the analytical context their question begins in.

\section{Limitations}
\label{sec:conclusion:limitations}

The work has four limitations that should be acknowledged.

\paragraph{Single-day GPS sample.}
The local GPS dataset covers a single day (2025-03-18) over
nine routes. The full fifty-day dataset is referenced in the
dataset specification~\cite{hcmcgps2025} and is reachable but
not loaded locally. Seven charts in the catalogue
(Section~\ref{sec:proposed:milestones}) remain at SHOULD or
WON'T-tier until the full dataset is loaded; the activation is
scheduled at the start of P9
(Section~\ref{sec:future:gps}).

\paragraph{Calendar mismatch between GPS and AFC.}
GPS data spans March 2025 while AFC tickets span
January--March 2026. The two streams cannot be joined by
timestamp and are treated as complementary sources whose
findings are combined narratively in the dashboard at the
\texttt{route}~$\times$~\texttt{hour-of-day}~$\times$~\texttt{day-type}
aggregate grain (Section~\ref{sec:proposed:schema}). A future
data release that produces overlapping GPS and AFC windows
would let the cross-source analysis happen at finer grains.

\paragraph{No alighting data.}
AFC records boarding events only. ``Passengers on the bus at
time~$T$'' is therefore not a directly computable quantity;
all crowding analyses are framed in terms of boardings per
stop or per trip. The limitation is documented visibly in the
Stop Demand Explorer (Chart~3.3) and the Network Health card
so end users encounter it before forming unrealistic
expectations.

\paragraph{QoS-2 sensitivity to GPS noise.}
In the dense urban canyons of HCMC, GPS positioning noise
causes pings to fall outside the 50-metre route buffer even
when the bus is in fact on-route. This systematically lowers
absolute QoS-2 scores; the metric is therefore safer to
interpret relatively (comparing routes to each other) than
absolutely. The Reliability scorecard (Chart~5.4) will surface
this caveat as an informational callout next to the QoS-2 KPI.

\section{Contributions}
\label{sec:conclusion:contributions}

Per the advisor's instruction, the contributions of this work
are summarised at the end of the report so that the reader
encounters them after the evidence rather than before.
Phase~1 of this thesis focuses on solutioning and prototyping;
the contributions are therefore deliberately scoped and short.

\paragraph{1. A reconciled design specification.}
Two parallel data streams (GPS and AFC) and four stakeholder
groups (passenger, driver, company, MCPT operator) are
reconciled into a single dashboard specification of six themed
cards and thirty-one canonical charts, with explicit invariants
for linked brushing, story-subtitled charts, and a universal
entity detail panel
(Chapters~\ref{chapter:proposed} and~\ref{chapter:prototype}).

\paragraph{2. A four-layer warehouse design.}
A four-layer DuckDB warehouse design (raw views, conformed
dimensions, fact views, pre-computed aggregate tables) is
specified, partially scaffolded (Layer~1), and justified
against two design alternatives. Each of the thirty-one charts
in the design is mapped to exactly one of eight aggregate views
as a hard data contract
(Sections~\ref{sec:proposed:warehouse_local}
and~\ref{sec:proposed:modules}).

\paragraph{3. An evidence-grounded module scoping.}
Four executed Jupyter notebooks (\texttt{notebooks/01--04})
produce the eleven headline EDA findings on which the SDVF
tier prioritisation is grounded, including the thirty-four
percent student-fare share that anchors the flagship card and
the calendar mismatch that forces a row-level join between GPS
and AFC to be ruled out
(Chapter~\ref{chapter:prototype}).

\paragraph{4. A high-fidelity, role-anchored UI mockup.}
A pixel-accurate light-theme mockup of the hub, one card
interior, and the universal entity-detail panel is published
at \texttt{docs/design/app-layouts.html}, locking the visual
language before any Streamlit code is written
(Section~\ref{sec:prototype:mockups}).

\paragraph{Mapping to the capstone topic deliverables.}
The capstone topic description~\cite{capstoneDescription}
names five Phase~1 tasks. Their coverage in this report is:
\emph{(i)~studying the HCMC bus dataset} →
Section~\ref{sec:prototype:eda};
\emph{(ii)~proposing analytical problems} →
Sections~\ref{sec:proposed:problem}
and~\ref{sec:proposed:modules};
\emph{(iii)~designing the analysis and visualisation
solution} → Chapters~\ref{chapter:proposed}
and~\ref{chapter:prototype};
\emph{(iv)~researching and selecting techniques and tools} →
Chapter~\ref{chapter:background} and the alternatives-considered
subsections~\ref{sec:proposed:warehouse_alts},
\ref{sec:proposed:module_rationale},
\ref{sec:proposed:dashboard_alts};
\emph{(v)~developing, testing, and conducting initial
evaluation} → partially met by the warehouse skeleton
(Section~\ref{sec:prototype:warehouse}) and the executed
notebooks (Section~\ref{sec:prototype:eda}); full Phase~2
development continues per the roadmap of
Chapter~\ref{chapter:future}.

A formal implementation and user study are out of scope for
Phase~1 and are sketched in Chapter~\ref{chapter:future}.
